\chapter*{Resumen\markboth{Resumen}{Resumen}}
\addcontentsline{toc}{chapter}{Resumen}

\renewcommand{\keywords}{\textbf{\emph{Palabras Clave ---~}}}

El presente informe documenta el desarrollo de \textit{DOMU}, una plataforma integral para la administración y gestión de comunidades residenciales que combina una pagina web, una Progressive Web App (PWA) móvil y un sistema de control de acceso mediante pistola escáner QR. El proyecto surge a partir de la necesidad de digitalizar procesos manuales ---registro de visitas, trazabilidad de encomiendas, cobranza de gastos comunes, gestion de espacios, comunicación interna, entre otros--- identificada en comunidades chilenas en el contexto de las exigencias de la Ley de Copropiedad Inmobiliaria N\textsuperscript{o}~21.442.

\medskip
\noindent\textbf{Objetivo general}. Desarrollar una plataforma web integral que integre y optimice los procesos criticos de la administracion de copropiedades: registro y verificacion de visitas con pre-autorizacion y apoyo de escaneo del QR de la cedula de identidad chilena, trazabilidad de encomiendas, gestion de gastos comunes con generacion de estados de cuenta y validacion manual de pagos con comprobante, gestion de personal operativo y tareas, reporte de incidentes, reserva de areas comunes y comunicacion interna mediante muro comunitario, chat y marketplace.

\medskip
\noindent\textbf{Metodología}. Se adoptó una \textit{metodología híbrida} que combina la estructura secuencial del modelo en Cascada con la flexibilidad del desarrollo Iterativo e Incremental. Las fases de análisis y diseño siguieron un enfoque secuencial con documentación formal, mientras que la construcción se realizó en ciclos cortos donde se implementaron, integraron y probaron módulos funcionales progresivamente.

% \medskip
% \noindent\textbf{Estructura del informe}. El documento se organiza en nueve capítulos:
% \begin{itemize}
%   \item \textbf{Capítulo~1 -- Introducción}: contexto, problemática, procesos BPMN, propuesta de solución, análisis de competidores, justificación y objetivos.
%   \item \textbf{Capítulo~2 -- Proyecto}: objetivos del software, metodología, tecnologías, estándares y herramientas.
%   \item \textbf{Capítulo~3 -- Especificación de Requerimientos}: límites, restricciones, 20~requerimientos funcionales, requerimientos no funcionales e interfaces externas.
%   \item \textbf{Capítulo~4 -- Factibilidad}: análisis técnico, operativo y económico con proyección a 5~años y cálculo del VAN.
%   \item \textbf{Capítulo~5 -- Análisis}: historias de usuario, actores, 27~casos de uso y matriz de trazabilidad RF--CU.
%   \item \textbf{Capítulo~6 -- Diseño}: servicios web, modelo de datos, diagramas entidad-relación, guía de estilo, mockups y arquitectura del software.
%   \item \textbf{Capítulo~7 -- Desarrollo}: estructura del código (190~clases backend, 92~archivos frontend), fases de implementación y detalles técnicos (QR, JWT, WebSocket, PDF, Flyway).
%   \item \textbf{Capítulo~8 -- Implantación}: arquitectura de despliegue basada en contenedores, plan de capacitación y estrategia de lanzamiento.
%   \item \textbf{Capítulo~9 -- Conclusión}: cumplimiento de objetivos, competencias, dificultades y trabajo futuro.
% \end{itemize}

\medskip
\noindent\textbf{Resultados}. El sistema alcanzo un estado funcional completo con mas de 130~endpoints REST, 2~canales WebSocket, 35~migraciones de base de datos, 54~vistas de interfaz y una PWA instalable en dispositivos moviles. El control de acceso mediante pistola escaner QR permite registrar visitas en menos de cinco segundos a partir de la lectura del codigo QR de la cedula de identidad. El modulo financiero genera estados de cuenta en PDF, permite la carga de comprobantes por parte del residente y soporta la validacion manual del pago por administracion. El analisis de factibilidad economica arroja un VAN positivo de \$158.219.839~CLP al 10\%, confirmando la viabilidad comercial del proyecto.

\medskip
\noindent\textbf{Conclusiones}. Los seis objetivos especificos fueron cumplidos, logrando una plataforma que centraliza la gestion residencial en un solo sistema. Cuatro casos de uso fueron postergados para una version futura del producto: gestion de morosidad y conciliacion, mantenimiento preventivo, administracion de estacionamientos y bitacora tecnica de ascensores. Se concluye que DOMU constituye una propuesta competitiva y viable para el mercado chileno de proptech.

\keywords{administración de condominios, gestión de comunidades residenciales, plataforma web, control de acceso QR, gastos comunes, proptech}
